% Dilution d'un facteur 10 puis mesure du pH (L'essentiel, chapitre 1 ; exercice 15)
\begin{tikzpicture}[x=1cm,y=1cm,
  etape/.style={font=\otLight\fontsize{7.5}{9.2}\selectfont,text=otInk,align=center,text width=31mm,anchor=north}]
% 1. Prélèvement à la pipette jaugée
\pic at (0,0) {becher={largeur=1.5,hauteur=1.8,niveau=1.0}};
\pic at (0.05,0.25) {pipette={hauteur=4.1,niveau=1}};
\node[pcLegendeG] at (-0.45,0.55) {S$_0$};
\node[etape] at (0,-0.35) {\otRond{1}\enspace Prélever 10,0 mL de S$_0$ à la pipette jaugée, munie d'une propipette.};
% 2. Transfert dans la fiole jaugée
\pic at (3.5,0) {fiole={hauteur=3.1,niveau=0.18}};
\pic at (3.5,2.55) {pipette={hauteur=2.9,niveau=0.35}};
\draw[pcRepere] (3.2,1.9) -- (2.7,2.25) node[pcLegende,anchor=east]{100,0 mL};
\node[etape] at (3.5,-0.35) {\otRond{2}\enspace Verser dans une fiole jaugée de 100,0 mL.};
% 3. Compléter au trait de jauge
\pic at (7.1,0) {fiole={hauteur=3.1,niveau=1}};
\pic at (5.85,1.25) {pissette};
\draw[pcRepere] (7.34,2.35) -- (7.85,2.75) node[pcLegende,anchor=west]{trait de jauge};
\node[pcLegendeG] at (7.1,0.55) {S$_1$};
\node[etape] at (6.95,-0.35) {\otRond{3}\enspace Compléter avec de l'eau distillée jusqu'au trait de jauge, boucher, agiter.};
% 4. Mesure du pH
\pic at (10.6,0) {agitateur={largeur=2.2}};
\pic at (10.4,0.5) {becher={largeur=1.4,hauteur=1.7,niveau=0.95}};
\pic at (10.25,0.58) {barreau};
\pic (sd) at (10.65,0.72) {sonde={hauteur=2.3}};
\pic (af) at (12.55,1.55) {afficheur={largeur=1.5,hauteur=1.1,texte={2,04}}};
\draw[pcFil] (sd-haut) to[out=90,in=180] (af-prise);
\node[etape] at (11.0,-0.35) {\otRond{4}\enspace Mesurer le pH de S$_1$ avec un pH-mètre étalonné.};
\end{tikzpicture}
\par\smallskip\legende{Préparation d'une solution fille S$_1$ dix fois moins concentrée que S$_0$, puis mesure de son pH. Schéma sans échelle.}
